\section{Model}

A firm's production consists of a sequence of steps $\mathcal{S}=(s_1,\ldots,s_m)$.\footnote{
We hold the set of production steps fixed and abstract from the creation or disappearance of steps in our model.
}
Each step can be performed manually, performed with AI assistance (augmentation), or, when chained with neighboring steps, executed by AI with no direct human involvement (automation).
Completing a step requires worker time and skill, and these requirements differ across the modes of execution.
The firm groups contiguous subsequences of steps into \emph{tasks}, choosing how to deploy AI across the workflow to minimize production costs.
Figure~\ref{fig:steps_tasks} illustrates this procedure.
We formally characterize the components of the figure below in turn, and then introduce the firm's optimal AI deployment strategy and show that it can be computed efficiently in polynomial time.

\begin{figure}[ht!]
  \begin{center}
  \caption{Bundling Steps into Tasks via AI Strategy}
  \label{fig:steps_tasks}
  \vspace{0.2cm}
  % Steps-to-Tasks Visualization.
% This is the top two layers of plots/TikZ_visualization/hierarchical_production_illustration.tex,
% cut at the task layer and kept identical to it in scale, node distance, box
% styles, and colors, so that a step, a task, and an aggregation arrow read the
% same in both figures. The only omissions are the job layer and the red snaked
% hand-off arrows, neither of which exists before jobs are introduced.
\begin{adjustbox}{center,trim=0pt 0pt 0pt 0pt,clip,margin=-2cm 0cm}
  \begin{tikzpicture}[scale=0.73, transform shape,
    node distance=0.35cm and 0.91cm,
    every node/.style={rectangle, rounded corners, draw, align=center, minimum width=2cm, minimum height=1cm},
    manual/.style={fill=gray!10},
    automated/.style={fill=green!30},
    augmented/.style={fill=orange!30},
    dashedbox/.style={draw, rectangle, rounded corners, dashed, inner sep=0.25cm, line width=1.25pt},
    humanbox/.style={rectangle, rounded corners, draw, align=center, minimum width=2.25cm, minimum height=1cm, fill=olive!20},
    aibox/.style={rectangle, rounded corners, draw, align=center, minimum width=2.25cm, minimum height=1cm, fill=blue!20},
    >=latex
  ]

    % Top layer (Steps)
    \node[manual] (S1) {Step 1\\(Manual)};
    \node[augmented, right=of S1] (S2) {Step 2\\(Augmented)};
    \node[manual, right=of S2] (S3) {Step 3\\(Manual)};
    \node[automated, right=of S3] (S4) {Step 4\\(Automated)};
    \node[automated, right=of S4] (S5) {Step 5\\(Automated)};
    \node[augmented, right=of S5] (S6) {Step 6\\(Augmented)};
    \node[manual, right=of S6] (S7) {Step 7\\(Manual)};

    % Arrows between steps
    \draw[->, thick] (S1) -- (S2);
    \draw[->, thick] (S2) -- (S3);
    \draw[->, thick] (S3) -- (S4);
    \draw[->, thick] (S4) -- (S5);
    \draw[->, thick] (S5) -- (S6);
    \draw[->, thick] (S6) -- (S7);

    % Dashed boxes grouping consecutive steps into tasks
    \node[dashedbox, draw=olive, fit=(S1)] (DS1) {};
    \node[dashedbox, draw=blue, fit=(S2)] (DS2) {};
    \node[dashedbox, draw=olive, fit=(S3)] (DS3) {};
    \node[dashedbox, draw=blue, fit=(S4)(S5)(S6)] (DS4-6) {};
    \node[dashedbox, draw=olive, fit=(S7)] (DS7) {};

    \path (S1.west) -- (S7.east) coordinate[midway] (CenterTop);

    % Bottom layer (Tasks)
    \node[humanbox, below=2.4cm of CenterTop, xshift=-6cm] (T1) {Task 1\\(Human)};
    \node[aibox, right=of T1] (T2) {Task 2\\(AI Chain)};
    \node[humanbox, right=of T2] (T3) {Task 3\\(Human)};
    \node[aibox, right=of T3] (T4) {Task 4\\(AI Chain)};
    \node[humanbox, right=of T4] (T5) {Task 5\\(Human)};

    % Arrows between tasks
    \draw[->, thick] (T1) -- (T2);
    \draw[->, thick] (T2) -- (T3);
    \draw[->, thick] (T3) -- (T4);
    \draw[->, thick] (T4) -- (T5);

    % Aggregation arrows, colored to match the dashed box each task comes from
    \draw[->, thick, olive] (DS1.south) -- (T1.north);
    \draw[->, thick, blue]  (DS2.south) -- (T2.north);
    \draw[->, thick, olive] (DS3.south) -- (T3.north);
    \draw[->, thick, blue]  (DS4-6.south) -- (T4.north);
    \draw[->, thick, olive] (DS7.south) -- (T5.north);

  \end{tikzpicture}
\end{adjustbox}

  \end{center}
  \vspace{0.1cm}
  \footnotesize{\emph{Notes:} The top layer shows seven production steps, each completed in manual (gray), automated (green), or augmented (orange) mode; dashed boxes group contiguous steps into tasks. The bottom layer shows the resulting tasks: a run of automated steps ending in an augmented step forms an AI chain (blue; here steps 2--4 become Task~2), and a lone augmented step is an AI chain of length one (here step~6 becomes Task~4), while a step performed manually on its own is a human task (olive). Black arrows show the production flow within each layer and the grouping of step blocks into tasks.}
\end{figure}



\subsection{Steps}
\label{sec:model.steps}

Steps are the smallest unit of work in our framework.
An individual step can be completed either with AI or manually (without AI).
If the firm assigns step $i$ to manual execution (denoted by $\manualLetter$), it hires a human worker of skill level $\manualSkill{i}$ and pays for the time $\manualTime{i}$ required to complete the step without AI.
Alternatively, if the firm assigns step $i$ to AI-assisted execution (denoted by $\AIletter$), it hires a human worker of (potentially different) skill level $\AIskill{i}$ and pays for the time $\AItime{i}$ spent completing the step in collaboration with AI.\footnote{
This could, for example, represent the time spent formulating a prompt for AI and checking the response, or ``managing'' the AI.
}
Such a collaboration succeeds independently with some (step-dependent) probability $q_i=\alpha^{d_i}$, where $\alpha$ is the quality of the general purpose AI technology and $d_i$ represents how ``difficult'' step $i$ is for AI.
The collaboration must be repeated until it succeeds, incurring the time cost of $\AItime{i}$ per iteration for a total expected time cost of $\AItime{i}/q_i$.

Definitions~\ref{def:manual_task} and~\ref{def:augmented_task} formalize the two modes of performing a single step.

\begin{definition}[Manual Step]
\label{def:manual_task}
A step $i$ is performed manually if it is executed entirely by a human worker without AI assistance.
The associated skill and time costs for a manual step are denoted by $(\manualSkill{i}, \manualTime{i})$.
\end{definition}

\begin{definition}[Augmented Step]
\label{def:augmented_task}
A step $i$ is augmented if it is executed by the AI, after which its output is reviewed and approved by a human worker.
The associated skill and (expected) time costs of managing a single augmented step are denoted by $(\AIskill{i}, \frac{\AItime{i}}{q_i})$, where $q_i = \alpha^{d_i} \in (0,1]$ is the AI's probability of successfully completing the step, increasing in the quality $\alpha$ of the general-purpose AI technology and decreasing in the step's difficulty $d_i$ for AI.
\end{definition}



\subsection{Tasks}
\label{sec:model.tasks}

Tasks are the smallest unit of work for a human worker.
A step performed on its own, whether manually or with AI augmentation, constitutes a task, with skill and (expected) time costs equal to those of the step in its chosen mode: $(\manualSkill{i},\, \manualTime{i})$ if performed manually and $(\AIskill{i},\, \AItime{i}/q_i)$ if augmented.

Steps can also be executed in an \emph{automated} mode in which AI performs the step end-to-end with no direct human involvement, and hence at zero direct human cost.\footnote{That execution of automated steps do not incur any cost to the firm is consistent with the AI subscription business model, where the firm pays an upfront fixed cost to an AI provider and uses AI tools at zero marginal cost afterward. The mechanisms discussed in the model remain unaffected even if we assume each time the firm deploys AI on a step it has to incur a marginal cost, but we abstract away from considering this to avoid complicating the model.}
Automation, however, arises only within an \emph{AI chain}: a contiguous run of AI-executed steps in which each step's output feeds directly into the next and a human verifies only the final output.
Every step of such a chain but the last is automated while the last step is augmented, since its output is the one the worker reviews.
The AI chain as a whole can be thought as a single \emph{composite} task delegated by the human to AI for execution.

Consider an AI chain spanning steps $(s_\ell, \dotsc, s_r)$.
The worker manages the chain by prompting for and evaluating only its final step $s_r$, since all earlier steps run automatically ``under the hood'' and remain beneath the worker's awareness.
Managing one attempt at the chain therefore costs the same as augmenting step $s_r$: skill $\AIskill{r}$ and time $\AItime{r}$ per attempt.
An attempt succeeds only if AI completes every step of the chain, which, assuming independent failures, occurs with probability $\prod_{i=\ell}^{r} q_i$; the chain is retried until it succeeds, for a total expected time of $\AItime{r}/\prod_{i=\ell}^{r} q_i$.
The chain's skill and (expected) time costs are therefore $(\AIskill{r},\, \AItime{r}/\prod_{i=\ell}^{r} q_i)$.
Since $q_i = \alpha^{d_i}$, the chain succeeds with probability $\alpha^{\sum_{i=\ell}^{r} d_i}$, so $\sum_{i=\ell}^{r} d_i$ is the chain's total difficulty, aggregating additively over its constituent steps.
Finally, an AI chain may consist of a single step, with $\ell = r$: having no automated predecessors, it is simply an augmented step.
An AI-augmented step is thus an AI chain of length one, a convention we adopt throughout the paper.

Definitions~\ref{def:automated_task} and~\ref{def:ai_chain} formalize the automated execution mode and the AI chain.

\begin{definition}[Automated Step]
\label{def:automated_task}
A step is automated if it is executed entirely by AI without direct human intervention, and its output is passed directly to a subsequent augmented or automated step.
The direct human costs (in both the skill and time dimension) of an automated step are zero.
\end{definition}

\begin{definition}[AI Chain]
\label{def:ai_chain}
An AI chain is a contiguous block of one or more sequential steps executed by AI, in which all steps but the final one are automated and the final step is augmented.
An AI chain spanning steps $(s_\ell,\dots,s_r)$ has steps $(s_\ell,\dots,s_{r-1})$ automated and step $s_r$ augmented, with skill and (expected) time costs $\left(\AIskill{r},\, \AItime{r}/\prod_{i=\ell}^{r} q_i\right)$.
\end{definition}

Finally, we define an \emph{AI deployment strategy} (or \emph{AI strategy} for short) as a partition of the step sequence $\mathcal{S}$ into a sequence of contiguous tasks $\mathcal{T}$, each of which is either a manual step or an AI chain.
The strategy in Figure~\ref{fig:steps_tasks}, for instance, groups seven steps into five tasks: two AI chains (a three-step chain over Steps 2--4 and a length-one chain at Step~6) and three manual steps.
Table~\ref{tab:model_parameters} summarizes the notation introduced for steps and tasks.

\begin{table}[ht!]
\centering
\caption{Notation Summary}
\label{tab:model_parameters}
\begin{tabularx}{\textwidth}{l X}
\toprule
\textbf{Symbol} & \textbf{Description} \\
\midrule
$\alpha \in (0,1]$ & Quality of the general-purpose AI technology \\
\addlinespace
$s_i$ & Production step $i$ \\
$\mathcal{S} = (s_1, \dotsc, s_m)$ & Sequence of the firm's $m$ production steps \\
$\manualSkill{i},\ \manualTime{i}$ & Skill and time cost of completing step $i$ manually \\
$\AIskill{i},\ \AItime{i}$ & Skill and time cost of managing one AI attempt at step $i$ \\
$d_i$ & Difficulty of step $i$ for AI \\
$q_i = \alpha^{d_i}$ & Probability that AI completes step $i$ successfully \\
\addlinespace
$T_b$ & Task $b$: a manual step or an AI chain (with a standalone augmented step considered a length-one AI chain) \\
$\mathcal{T} = (T_1, \dotsc, T_n)$ & AI deployment strategy; a partition of the steps into $n$ tasks \\
$\skillcost{b},\ \timecost{b}$ & Skill and (expected) time cost of task $b$: $(\manualSkill{i},\, \manualTime{i})$ if a manual step, $(\AIskill{r},\, \AItime{r}/\prod_{i=\ell}^{r} q_i)$ if an AI chain over steps $\ell, \dotsc, r$ \\
\bottomrule
\end{tabularx}
\end{table}



\subsection{Jobs and Wages}
\label{sec:jobs_wages}

A \emph{job} is the set of tasks assigned to a single worker, that is, a contiguous block of the task sequence $\T$.
The firm can in principle split its workflow into several jobs and hire a separate worker for each, and it is free to choose where one job ends and the next begins.
This choice trades the gains from specializing workers on narrower sets of tasks against the coordination costs of splitting work across more of them, and it is what reallocates tasks across workers.
We delegate the full characterization of this job-design problem, including what pins down the boundary of a job, to Appendix~\ref{app:jobdesign}.
In the body we instead hold the boundary of a job fixed and exogenous, and treat the entire step sequence $\mathcal{S} = (s_1, \dotsc, s_m)$, which the AI strategy organizes into the task sequence $\T$, as a single worker's job.
This job is the unit whose cost the firm minimizes.

A worker is paid for the entire time their job takes to complete, at a rate equal to the skill the job requires.\footnote{
One motivation for this formulation is that workers make a human-capital investment to acquire the skills their job demands, an investment that must be recovered through wages; we express skill levels directly in units of the wage they command.
\pscomment{Give a brief intuition here as the aggregation stuff don't exist anymore.} Rather than formalize such a wage model here, we impose it as an assumption and develop a more detailed formulation in Section~\ref{sec:aggregation}.
}
The skill a job requires is the total skill of its tasks, so the worker's wage rate is
\begin{align}
\label{eq:wage}
\text{Wage} = \sum_{T_b \in \T} \skillcost{b},
\end{align}
and the firm pays this rate for as long as the job takes, regardless of which task the worker is performing at any given moment.\footnote{
This implicitly assumes a common base wage rate, so that all wage differences are driven by the human-capital cost of acquiring skills.
More generally, manual and AI-assisted tasks could carry different base wage rates $w_{\manualLetter}$ and $w_{\AIletter}$, reflecting the supply of and demand for each type of labor; Equation~\eqref{eq:wage} normalizes these rates to $1$ and folds them into the skill costs $\skillcost{b}$.
We maintain this normalization for notational simplicity but it can be relaxed to more general cases.
\label{foot:wage}
}
The duration of a job is the total (expected) time requirement of its tasks,
\begin{align*}
\sum_{T_b \in \T} \timecost{b},
\end{align*}
The labor cost (i.e., wage bill) of a job is therefore its wage rate times the time it takes,
\begin{align*}
\text{Wage Bill} = \bigl(\sum_{T_b \in \T} \skillcost{b}\bigr)\bigl(\sum_{T_b \in \T} \timecost{b}\bigr).
\end{align*}


\subsection{The Firm's Problem}
\label{sec:model.jobs}

The firm deploys AI across its workflow to minimize production costs.
We assume labor is the only cost of production and thus the firm minimizes its wage bill paid to workers it hires.
What the firm can adjust to lower its labor cost depends on the time horizon of optimization, and we distinguish three.

In the \emph{short run}, the boundary of the job and wages are fixed as the firm has already hired a worker to do a designated set of tasks.
The wage is then sunk, since the firm owes it for as long as the job takes no matter how AI is deployed, so the only margin AI can move is the time the job takes.
The firm therefore chooses its AI deployment strategy $\T$ to minimize the total expected time of completing the workflow,
\begin{align}
\label{eq:totalcost_shortterm}
\min_{\T} \ \sum_{T_b \in \T} \timecost{b}.
\end{align}
AI is deployed on a step, or on a run of steps as a chain, whenever doing so lowers this expected time.
%This short-run problem is the one we solve in the remainder of the body.

In the \emph{medium run}, the boundary of the job remains fixed as plausibly reorganizing production and reassigning workers to tasks is implausible to take place, but free entry into jobs lets the firm adjust the wage it pays to the skill the job actually requires if AI lowers the cost executing the set of tasks in the job.
Deploying AI can then lower cost on two margins at once, since an augmented or automated step may demand both less skill and less time than manual execution.
Because the worker is paid in proportion to the job's skill for the entire time the job takes, the firm minimizes the product of the job's total skill and its total time,
\begin{align}
\label{eq:totalcost_mediumrun}
\min_{\T} \ \Bigl(\sum_{T_b \in \T} \skillcost{b}\Bigr)\Bigl(\sum_{T_b \in \T} \timecost{b}\Bigr).
\end{align}

In the \emph{long run}, the firm additionally redraws the boundaries of jobs, reallocating tasks across workers in addition to optimizing the AI strategy in production.
Narrowing a job lets the firm staff it with workers specializing to do specific tasks (thus requiring lower wages for fewer tasks), but splitting the workflow across more workers introduces coordination costs between them, and the firm chooses the AI deployment and the job design jointly to balance the two.
We develop this long-run problem, which nests the medium run as the special case of fixed job boundaries, in Appendix~\ref{app:jobdesign}.



\subsection{Optimal AI Deployment}
\label{sec:optimization.short}

We now turn to solving the firm's short-run problem in \eqref{eq:totalcost_shortterm}.
Although the number of possible AI strategies over the $m$ steps grows exponentially in $m$, the cost-minimizing strategy can be computed via dynamic programming in $O(m^2)$ time.

\begin{proposition}
\label{prop:shortrun_dp}
Given a workflow with $m$ steps, the cost-minimizing AI strategy can be calculated in time $O(m^2)$ via dynamic programming.
\end{proposition}

The proof, a dynamic program over prefixes of the step sequence, is given in Appendix~\ref{app:omitted_proofs}.
